\chapter*{How to Use This Book}
\addcontentsline{toc}{chapter}{How to Use This Book}

The book is designed to be read alongside its companion project. Each major
technique is approached in four connected forms:

\begin{enumerate}
    \item the chapter explains the concept, mathematics and algorithm;
    \item the companion notebook develops and explores the implementation;
    \item the finished implementation is added to the \pyimage{} package; and
    \item automated tests verify the behaviour of the reusable implementation.
\end{enumerate}

The normal progression is therefore:

\begin{center}
\textbf{theory} $\rightarrow$ \textbf{algorithm} $\rightarrow$
\textbf{notebook} $\rightarrow$ \textbf{library} $\rightarrow$
\textbf{tests} $\rightarrow$ \textbf{library comparison}.
\end{center}

The library comparison comes last deliberately. Where a technique can be
implemented directly, the book first constructs it from basic Python and/or
NumPy operations. Established libraries are then used to compare behaviour,
performance and API design.

\section*{Companion notebooks}

The Jupyter notebooks mirror the chapter numbering. For example,
\texttt{notebooks/17-edge-detection.ipynb} accompanies Chapter~17. Notebooks
may include exploratory code and intermediate implementations that do not
belong in the final package.

\section*{The \pyimage{} package}

The canonical reusable implementations live under
\texttt{project/src/pyimage/}. Once an implementation has been developed and
validated in a notebook, the polished function is added to the package and
covered by the test suite.

\section*{Historical sign-off items}

Some chapters contain a box labelled \emph{Original portfolio/sign-off item}.
The stars reproduce the relative weighting used in the historical material
where that weighting can be verified. They are retained as a record of the
original practical scope, not as assessment marks for this book.
